%=====================================================================
\chapter{Resources}\label{app:resources}
%=====================================================================
\normalfigures

\section{The Set Texts}

These are the books HEC names for this course. The first is the one to buy.

\begin{tight}
  \item S.\,J. Russell and P.\ Norvig, \emph{Artificial Intelligence: A Modern
        Approach}, 3rd edition, Prentice Hall, 2010. Named by both the HEC 2017 and
        HEC 2023 outlines. Chapters 2--7 and 13--14 of Russell and Norvig cover
        most of this handout; it is more thorough and much longer, and it is the
        right place to go when a chapter here leaves you wanting the general case.
  \item G.\,F. Luger and W.\,A. Stubblefield, \emph{AI Algorithms, Data
        Structures, and Idioms in Prolog, Lisp and Java}, Pearson Addison-Wesley,
        2009. Named by HEC 2017. Useful precisely because it implements the
        algorithms rather than describing them.
  \item R.\,O. Duda, P.\,E. Hart and D.\,G. Stork, \emph{Pattern Classification},
        Wiley, 2001. Named by HEC 2017. It is a pattern-recognition text and
        belongs to \emph{Machine Learning} rather than to this course; read it
        after that one.
  \item D.\ Poole and A.\ Mackworth, \emph{Artificial Intelligence: Foundations of
        Computational Agents}, 2nd edition, Cambridge, 2017. Named by the HEC 2023
        outline for \emph{Knowledge Representation and Reasoning}. Free online, and
        the clearest available treatment of the agent framing this course adopts.
\end{tight}

\section{Where Each Chapter Sits in Russell and Norvig}

\rowcolors{2}{white}{lightbg}
\begin{longtable}{@{}L{4.2cm}L{4.4cm}L{6.2cm}@{}}
\toprule
\rowcolor{primary!15}
\textbf{This handout} & \textbf{Russell \& Norvig 3rd ed.} & \textbf{Note} \\
\midrule
\endhead
1 What AI Is & Ch.\ 1 & The history here is compressed into five waves \\
2 Agents and PEAS & Ch.\ 2 & Very close; read both \\
3 Representations in Python & --- & No equivalent; this chapter is specific to
this course \\
4 Uninformed search & Ch.\ 3.1--3.4 & \\
5 Informed search & Ch.\ 3.5--3.6 & The optimality proof is theirs \\
6 Local search & Ch.\ 4.1--4.2 & \\
7 Constraint satisfaction & Ch.\ 6 & \\
8 Adversarial search & Ch.\ 5 & \\
9 Planning & Ch.\ 10 & They develop partial-order planning, which this course
only names \\
10 Knowledge representation & Ch.\ 12 & Their treatment is heavier on ontologies \\
11 Propositional logic & Ch.\ 7 & \\
12 First-order logic & Ch.\ 8--9 & \\
13 Expert systems & Ch.\ 9.3--9.4 & Rule systems are scattered in R\&N; Luger is
better here \\
14 Uncertainty and decisions & Ch.\ 13, 16 & \\
15 Bayesian networks & Ch.\ 14 & \\
16 Fuzzy logic & Ch.\ 14.7 (brief) & R\&N are sceptical of fuzzy logic and say so;
read their argument \\
17 Learning & Ch.\ 18.1--18.2 & This course stops where they accelerate \\
18 Natural language processing & Ch.\ 22--23 & Their treatment is largely
statistical; this chapter is the symbolic half \\
19 Classic systems & scattered & The primary sources below are better \\
20--21 Trends, responsibility & Ch.\ 26--27 & The 4th edition (2020) is
substantially better here \\
\bottomrule
\end{longtable}

\section{The Primary Sources Worth Reading}

Short, readable and historically decisive. Chapter~19 reconstructs the first four.

\begin{tight}
  \item A.\,M. Turing, ``Computing Machinery and Intelligence'', \emph{Mind},
        1950. The paper that proposed the test, and answered nine objections to it
        that are still the nine objections people raise.
  \item A.\ Newell and H.\,A. Simon, ``GPS, a Program that Simulates Human
        Thought'', 1961. Means--ends analysis, and the beginning of the cognitive
        modelling tradition.
  \item J.\ Weizenbaum, ``ELIZA --- A Computer Program for the Study of Natural
        Language Communication between Man and Machine'', \emph{CACM}, 1966. Read
        it with Weizenbaum's later dismay at how it was received.
  \item M.\ Minsky and S.\ Papert, \emph{Perceptrons}, MIT Press, 1969. The
        exclusive-or result of \chref{learning}, and the book usually blamed for
        the first neural winter.
  \item E.\,H. Shortliffe, \emph{Computer-Based Medical Consultations: MYCIN},
        1976. The expert system that established explanation as a requirement
        rather than a feature.
  \item J.\ Pearl, \emph{Probabilistic Reasoning in Intelligent Systems},
        Morgan Kaufmann, 1988. Bayesian networks, from the person who made them
        computable.
  \item J.\ Lighthill, ``Artificial Intelligence: A General Survey'', 1973. The
        report that caused the first winter. Worth reading because its criticism
        was largely correct.
\end{tight}

\section{Departmental Materials}

\begin{tight}
  \item \texttt{Machine Learning Fall 2026/} --- the course this one hands
        learning to. Its Chapter~13 continues \chref{learning} exactly where this
        book stops.
  \item \texttt{Data Science Fall 2026/} --- the end of the three-course sequence.
  \item \texttt{HEC AI-DS Curriculum Map/} --- the three HEC curriculum booklets
        and the department's analysis of them, including the verbatim outlines this
        course is written against.
  \item \texttt{AI\_ML\_DS\_Fall2026\_ProgressionAtAGlance.pdf} --- the one-page
        map of all three courses and who owns each shared topic. Appendix~F states
        the same boundaries in prose.
  \item \texttt{Summer\_2026/AI\_Handout.pdf} --- the eight-week revision summary.
        Useful for examination practice, not a substitute for this volume.
  \item \texttt{Summer\_2026/question\_papers/} --- past papers and a Moodle
        question bank for this course.
\end{tight}

\section{Tools}

Nothing beyond Python 3.10 and NumPy is required, and that is deliberate
(see the Course Specification). If you want to go further after the course:

\begin{tight}
  \item \textbf{SWI-Prolog} --- the natural home of Chapters 11--13. Writing
        \chref{expert}'s rule engine in Prolog takes about fifteen lines, and
        seeing that is worth an afternoon.
  \item \textbf{python-constraint} or \textbf{OR-Tools} --- production constraint
        solvers, for after you have written the one in \chref{csp} yourself.
  \item \textbf{pgmpy} --- Bayesian networks in Python, for after \chref{bayesnets}.
  \item \textbf{PDDL} and a planner such as \textbf{Fast Downward} --- the modern
        descendant of the STRIPS representation in \chref{planning}.
\end{tight}
